\documentclass[12pt]{book}

\usepackage{graphicx}
\usepackage{eso-pic}
\usepackage{tcolorbox}
\usepackage{polyglossia}
\setdefaultlanguage[indentfirst=false,forceheadingpunctuation=false]{russian}
\setotherlanguages{english}

\RequirePackage[a4paper, headsep = 0.5 \headsep, left=2.5cm, right=2.1cm, top=2cm, bottom=2.1cm]{geometry}

\setmainfont{URW Bookman}
\newfontfamily\cyrillicfont{URW Bookman}[Script=Cyrillic, Ligatures=TeX]
\setlength{\parindent}{20pt}

\usepackage{enumitem}
\setlist[itemize, 1]{label =\raisebox{-0.3\height}{\scalebox{1.8}{\textbullet}}}

\title{Феноменология чуждого}
\author{Бернхард Вальденфельс}
\date{2006}
\begin{document}

\AddToShipoutPictureBG*{\includegraphics[width=\paperwidth]{title.png}}

\null
    \thispagestyle{empty}%
    \addtocounter{page}{-1}%
    \newpage

\tableofcontents

\chapter{Грани чуждого}

Выделить чуждое как отдельную тему --- значит уже упустить его, начать со знакомого места и планировать к нему же вернуться. Обращение к чуждому наверняка приведет к переменам и даже катарсису, хотя в конечном счете привычность и возобладает и может даже углубиться. Чуждое не безобидно и может отдалить нас от самих себя --- потому ему стараются сопротивляться, избегать его или ассимилировать. Эти старания оставляют субъекта наедине с самим собой и требуют выстраивать и оборонять порядок, исключающий проявление чуждого изнутри. Чуждое может пробудить любопытство и рассказать нам о нас самих, но как только оно подходит к тайне свободы и разума --- срабатывает сигнал ''Хаос!'', и свобода и разум остаются в своих правах. Иначе им пришлось бы отказаться от самих себя. Чуждость неизбежно ведет к враждебности, обе стороны все больше убеждаются в своей правоте. Мы становимся ближе и бдительней друг к другу. Встроенные в эмпирическую сеть защитные механизмы исходят из \textit{своего} и стремятся закрепиться в \textit{общем}. Чуждое предполагается внешним, ожидается, что оно будет идентифицировать себя и выглядеть незваным гостем. Его тогда можно будет подвергнуть оценке и суждению и поместить на повседневный моральный, политический, религиозный, культурный и интеллектуальный карантин.

Если же чуждое воспринимать как неизбежно неопределенное, как всегда сокрытое в \textit{своем}, где оно тревожит, соблазняет или пугает, ускользает от понимания, то опыт чуждого, значит, воздействует на \textit{наш} опыт, становясь отчуждением от опыта. Чуждость самореферентна и заразительна, ее последствия предшествуют всякой ее тематизации. Феноменология чуждого, предвосхищенная Гуссерлем в нескольких его новаторских концепциях (чуждый опыт, чуждое эго, чуждое тело и чуждый мир) обладает задатками науки, могущей говорить о чуждом во всей его полноте, всерьез принимать его радикальность, не довольствуясь сведением к его пониманию и практическому распознаванию. Феноменолог, всерьез принимающий эти требования ксенологии, наверняка будет вдохновляться работами Марселя Мосса, Георга Зиммеля, Вальтера Беньямина, Людвига Витгенштейна или Михаила Бахтина, не говоря уже о таких литераторах как Кафка, Музиль, Целан, Валери и Кальвино. В мире искусства другие порядки, и художественное чутье легко преодолевает границы самоуверенности и принципиальности, если вступает в контакт с чуждым.

Настоящее исследование не пытается охватить проблематику чуждого во всей ее пугающей полноте, а лишь обозначает некоторые основные мотивы, побудившие автора разработать феноменологию чуждого в соответствующей серии книг. Первый том --- \textit{Жало чуждого} вышел в 1990г.\footnote{Формальное изложение феноменологии чуждого можно найти в Главе 1 \textit{Топографии чуждого} (\textit{Topographie des Fremden}, 1997). Аспекты чуждого рассматриваются в главах 5-6 книги \textit{Рваные линии опыта} (\textit{Bruchlinien der Erfahrung}, 2002).} Мотивы здесь развивают тему, а не просто ретроактивно обосновывают обращение к ней. Опыт чуждого --- как и опыт удивления, ужаса, любви, ненависти или тех едва заметных несостыковок, через которые возникнет новое --- нельзя вполне описать или обосновать. Разделение на главы выполнено так, что каждая рассматривает конкретный программный компонент проблематики чуждого. Они выстроены в порядке этапов формирования феноменологии чуждого --- они потому вводные, поясняющие тексты и дополняют уже опубликованные обширные исследования. В последующих абзацах я свяжу эту работу с предыдущими, что избавит меня от необходимости отдельно на них ссылаться и позволит сосредоточиться на ключевых темах: порядок -- претерпевание -- ответ -- тело -- внимание -- межкультурность.

В главе 1 рассматривается \textit{Порядок в сумерках} (\textit{Ordnung im Zwielicht}, 1987). На границах всякого порядка чуждое предстает экстраординарным, не находящим места в порядке --- и все же оно, будучи исключаемым, не ничто. Поскольку оно исключается не вообще --- а из конкретного порядка, то оно нечто большее, чем просто сливающееся с фоном в неопределенности. И поскольку оно взывает к нашим чувствам, то оно не просто шум, от которого мы хотели бы укрыться. Положение, в котором мы оказываемся, обнаруживая несхватываемое в схватываемом, беспорядочное в упорядоченном, невидимое в видимом, тишину в озвучиваемом, запускает цепь исследований, сводящихся к вопросу: как мы можем взаимодействовать с чуждым без предварительной нейтрализации или отрицания его эффектов, его вызовов и требований? Как феноменолог, я всегда старался избегать того, в чем нетерпеливая аудитория любит обвинять феноменологов чуждого --- возведения его в ранг некоего морального или религиозного авторитета или ''совершенно иного''. Это было бы обожествлением Другого --- в чем часто обвиняют Левинаса --- и морализацией чуждого. Но навязать переживание столь же возможно, как и, согласно Канту, чувство. Не то чтобы восприятие чуждого было лишь субъективным ощущением, как в том убеждает модерн, просто в этом мире вещей гораздо больше, чем может охватить школа мысли с ее нормами, ценностями и интересами.

К переживаниям можно относиться с разной степенью осторожности: позволить им говорить самим за себя, использовать как информационный ресурс или пытаться исключить --- будь то от небрежности, властолюбия или безразличия. В главах 2 и 3 эти возможности рассмотрены с точки зрения ответовой этики --- этики, возникающей в ответе, отклике, реакции, и так остающейся за рамками заповедей и запретов. До-предикатный опыт в гуссерлевском смысле здесь обращается в до-нормативный --- остающийся всегда позади, пока мы намечаем действия и строим планы, требующие моральной или правовой оценки. В главе 2 я буду ссылаться в основном на \textit{Рваные линии опыта} (\textit{Bruchlinien der Erfahrung}, 2002): анализ, проведенный в этом тексте, построен вокруг двойного события, состоящего из претерпевания и последующей реакции. В начале есть то, что предстает как события --- происходящим с нами, в чем я или Другой можем участвовать --- а не как волящий и действующий субъект. Ключевое здесь --- временной сдвиг. Не нечто иное предшествует инициативе ''я'' --- эта инициатива предшествует самой себе, не возникая из неоткуда. Эти связи предстанут в новом свете, когда в главе 3, рассматривая \textit{Регистр ответов} (\textit{Antwortregister}, 1994)\footnote{В данном случае я намеренно отклоняюсь от порядка, в котором были изложены эти идеи, чтобы облегчить понимания, поскольку \textit{Регистр ответов} предполагает некоторую часть того, что полностью раскрывается только в \textit{Рваных линиях опыта}. ''На что'' ответа указывает на ''чем'' преодоления.}, я раскрываю основные черты ответового поведения и возникающую в результате зова чуждого логику ответов. Парадокс творческого ответа представляет особый случай: оказывается, мы можем воспринимать ''что'' ответа как нечто изобретенное, но не можем сделать то же с его ''на что''. Зов чуждого может быть столь же слабо связан с этическим фундаментализмом, сколь и с технологическим конструктивизмом.

Глава 4 посвящена телесному опыту, и не в коем случае не просто расширяет сферу чуждого. В узле самореференции и самоустранения, определяющих его способом существования, тело оказывается как бы эмблемой чуждого. Чужеродность начинается не с меня самого --- это мало что значило бы,--- так называемый субъект выпадает из центра, когда предстает как \textit{Воплощенное Я} (\textit{Leibliches Selbst})\footnote{Название лекций о теле, опубликованых в 2000г., в которых эта тема рассматривается во всей ее глубине.}. В главе 5 уточняются и подкрепляются ссылками на литературу основные идеи \textit{Феноменологии внимания} (\textit{Phänomenologie der Aufmerksamkeit}, 2004). Внимание полностью подчинено двойному ритму становления очевидным и становления притягательным, двойное движению претерпевания и ответа потому преобретает форму, пронизывающую весь наш опыт и поведение. Что становится очевидным или мыслимым, что доходит до наших глаз или ушей, до наших органов чувств, что до ходит до нас --- не наше. Аллюзия названия \textit{Пороги смысла} (Sinnesschwellen, 1999) указывает, что вездесущее внимание ведет в область эстетики и особого этоса чувств. В последних двух главах говорится также о чуждых аспектах психоанализа и технологии, что занимало значительное место в \textit{Рваных линиях опыта} и \textit{Феноменологии внимания}. Как изначальный опыт пронизан моментами неизведанного и немотивированного, так и феномены в нем с самого начала отмечены подлинной феноменотехнологией и соматотехникой. Технология очень долго --- вопреки всем сетованиям современности об отчуждении в ней --- рассматривалась как расширение сферы собственности, но теперь все больше становится источником особой отчужденности.

Главу 6 можно читать как обзор, собрание различных аспектов чуждого в кратком, ориентированном на отношения между культурами, изложении. Межкультурность не ограничена философией.\footnote{Позвольте здесь упомянуть о выходящем в Вене журнале по межкультурной философии Polylog, в котором к настоящему времени появилось довольно много тематических выпусков.} Глобальное мышление не следует ни ожидать, ни желать --- оно неизбежно сформировало бы лишь новые гегемонистские отношения, разве что другими средствами. Философия, переставшая вращаться вокруг себя самой и пустившая корни в привычной культуре, дала бы облегчение. Попытка пересечь границы, не уничтожив их --- одно из будущих начинаний в области межкультурного взаимодействия, и оно значит нечто большее, чем просто сопоставление или смешение культур. Разнообразие слишком легко становится накоплением --- спокойным перебором любопытных вариаций. Опасность для культур, достигних зрелости, заключается в обращении к обществу охотников и собирателей.

Настоящая работа основана на серии отдельных исследований, появившихся в связи с разными обстоятельствами. Все эти исследования были переработаны и ориентированы на общую тему.\footnote{Эти идеи рассматриваются в \textit{Границы нормализации} (\textit{Grenzen der Normalisierung}, 1998).} Я постарался избежать ретроактивной систематизации и сохранить эвристическую множественность подходов, в то же время создав рельефы, благодаря которым --- как на объемной открытке --- были бы видны все значимые аспекты. Как писатели, мы иногда сожалеем, что письмо ограничено двухмерной плоскостью, но, к счастью, в письменной речи --- как и в чтении --- есть свои способы расставления акцентов, и эти подъемы и спуски не дают определенному направлению письма свестись к написанному.

\chapter{Человек как пограничное существо}

Чуждое --- это предельное как оно есть. Даже возникая в знакомом, оно приходит извне. Не бывает чуждого без чуждого места, и порядок нашей жизни, опыта, языка и действий определяет, сколь значимо для нас будет чуждое. Преобразование порядка значит и преобразование чуждого --- такого же многогранного, как и соответствующий порядок. О чуждом говорят не меньше, чем об эго --- пограничные зоны между и за порядками питают чуждое.

\section{Порядки и их пределы}

Обычно верно полагать, что не только порядок ограничен, но и что эти границы возникают в процессе упорядочивания. Нечто --- то, что оно есть в силу его отделения от прочих вещей: камней, растений, животных, природных или искусственных объектов. Однако, в мире людей понятие границ вызывает особое беспокойство, потому что границы всегда под сомнением. Поведение человеческого существа не определяется ни инстинктами, ни искусственными программами: оно не ограничено, а, скорее, особым образом соотнесено с границами. Это касается и места и времени --- границ, определяющих что такое здесь и сейчас,--- и запретов, сдерживающих желания и поступки, и пределов понимания, сдерживающих мышление. Неудивительно, что вопрос о границах существования остается в числе главных тем на протяжении всей человеческой истории: Иегова отделяет свет от тьмы, границы души неизмеримы, философ намечает границы чистого разума, системный теоретик воплощает возвышенный жест созидания в простом ''Проведи границу!''.

Видно, что установление границ, приводя к разным порядкам и структурам, носит не только прагматический и региональный, но и эпохальный характер. Каждая эпоха --- или каждая культура, общество, среда или форма жизни --- ведет себя в рамках некоторых границ, и отношение к этим границам, всегда сопровождаясь определенной политикой, значительно варьируется. Обращение с границами прямо показывает глубинный дух эпохи и может также послужить комментарием к тому, что так долго продвигало современность вперед или предшествовало ей, или подрывало ее и выходило за ее пределы. Рубеж к 21-ому столетию --- определенно очень специфическая граница, и у нас нет языка, подходящего для описания того, что ждет нас впереди.

\section{Безграничная вселенная}

Для контраста можно начать с безграничной вселенной. В нашей культурной традиции она наиболее емко представлена греческим космосом. Космос долгое время имел парадигмальное значение, олицетворяя классическую форму порядка. Он воплощал не один из множества порядков, а порядок вообще, как таковой, и потому единственной альтернативой ему были неупорядоченное многообразие хаоса. Этот космос придает каждому существу его собственную форму (πέρας), границы которой отделяют существо от себя и среды. Точная форма выражается в концептуальном определении (όρισμός); со времен Платона диалектика связывала каждое бытие как единое с прочим в цепочке отношений. Эта горизонтальная связь дополняется вертикальной иерархией, пропорции которой соответствуют степени отражения рационального целого в конкретном существе. Люди здесь стоят выше животных, греки --- выше варваров, мужчина --- выше женщины, созерцание --- выше действия. Рациональность, раскрывая закон целого, имеет решающее значение для места человека в иерархии существ.

Эта реляционная структура --- все различия в ней относительны --- имеет нижнюю и верхнюю границы. Нижняя формируется \textit{ин-дивидом}\footnote{in-dividual; англ. divide --- разделять.}, άτομον είδος --- он уже не может быть далее разделен без разрушения его идентичности. Верхняя граница --- Вселенная, не упорядоченная в рамках чего-либо другого; мир --- сущность, из которой нечего не исключено, постановит Аристотель в своей \textit{Физике} (2.6.207a8). Проще говоря, космос --- порядок без внешнего, с только лишь внутренними границами. И пересекающий эти границы либо попадает в дурную бесконечность бездонного и бессмысленного апейрона, либо, подобно Икару, поднявшемуся на запредельные высоты, терпит фатальную неудачу.

Но замкнутый всеобъемлющий космос, в котором есть место каждому, предполагает, что место проявления и выражения целого как \textit{целого} находится внутри самого этого целого. Психе --- согласно Аристотелю, ''в некотором смысле --- все'' --- становится ступенью порядка, к которому мимически приспосабливается. Космос потому предстает порядком, раскрывающим и выражающим \textit{себя как себя} через превращение всякого условия в момент себя. Такой порядок без внешнего соответствует \textit{мысли о внутреннем} --- \textit{penser du dedans}, перефразируя название работы Фуко\footnote{La Pensée du dehors, 1966} --- и в целом такая мысль была бы мыслью самой по себе. Однако даже древние греки знали о выпадающих из этой схемы фигурах: Сократе --- άτοπος, живущем в городе подобно прочим гражданам, но действующем как вопросительный знак, или о Платоновской \textit{мании} поэзии, или о софистах, выступающих против истинного логоса лексическими уловками, или об Эдипе, прозревшем только став слепым, прикованном к своей зловещей судьбе, изолировавшей его и заклеймившей бездомным (άπολις).\footnote{Фигура \textit{аполиса} вырисовывается на фоне полиса, и Софокл сравнивает его с ύψίπολις, ''самым выдающимся в городе'' (\textit{Антигона} 5.370). Человек-\textit{аполис} соответствует ''лицу без гражданства'' в нынешнем понимании, но в нем также виден аналог \textit{атопоса}: ''лишенный места'', ''неуместный'', ''странный''. Alpha privativum в обоих случаях можно понимать как не только отрицание, но и выражение отказа, отстранения, выхода.} Такие маргинальные фигуры встречаются и в других местах: приверженность текстам в общепринятой части мистического учения обычно сопровождается ответвлениями, отбрасывающими тень неортодоксальности и анархии --- это касается и еврейской и исламской, и христианской традиций. Даже Святой Павел отмежевывается от риторики, в которой γλώσσα или \textit{lingua} расслабляют язык до такой степени, что позволяют ему ускользнуть от индивидуального или общественного контроля. Высказывание, растворяющееся в космосе сказанного --- заменяется ничего не говорящим и потому бессмысленным. Обаяние, противопоставляемое Максом Вебером повседневной жизни институтов, пронизывает все и показывает, что любая нормальность --- даже космологическая, теологическая или космополитически моделируемая --- упускает нечто, находящее затем выражение в некой безумной бахроме аномалий.

\section{Принадлежность, чуждость, контингентность}

Глобальный порядок немыслим без разрыва, если только место, из которого разворачивается целое, не превращается в конкретное целое. Модерн --- движение, ставящее под сомнение видение целого. Оно пресекает самообман космотеороса, считающего себя частью созерцаемой им сцены, и мнимый порядок оказывается результатом процесса установления порядка [\textit{Ordnungsstiftung}]. Здесь два ключевых открытия: (1) \textit{я}, говорящее ''я'' до того, как будет названо ''субъект'' и взрыв структуры отношений в целом из-за этой самореферентности и (2) радикальная \textit{контингентность}, которая не только занимает открытые пространства в порядке, но и влияет на сам порядок. Такой порядок может не только выродиться в беспорядок, но и стать другим порядком, быть отличным от самого себя. В cogito Декарта есть не только вершина мысли, но и возможность создания Богом другой математики. Порядки, в которых мы движемся, оказываются лишь потенциальными. И странный протагонист \textit{Человека без свойств} Музиля подозревает, что даже Бог предпочитает говорить о мире лишь потенциально: ''ибо Бог создает мир и в то же время думает, что он мог бы быть другим'' (Музиль, 1961, 15) и связывает это с тезисом ''действительно любящий свою страну не считает ее лучшей''. Решающий фактор здесь --- что мотивы субъектности и рациональности выступают как двойной мотив. Субъект модерна больше не олицетворяет лишь рациональность и ищет свое место, которого нет.

Если всерьез --- во всей его радикальности --- принять этот двойной мотив, то проблема проведения границ изменится: \textit{я}, живущее в своей собственной сфере и своему собственному культурному порядку соответствующее даже в языке и чувствах, уже не будет сводимо к чему-то всегда одному и отдельному от другого в контексте универсальности. Самость и самопринадлежност --- результат проведения границ, разделяющих внутреннее от внешнего и так принимающих формы \textit{включения} и \textit{исключения}. Чувство принадлежности возникает, когда что-то выходит из него, и это уходящее мы принимаем чуждым и неоднородным. Такое разделение своего и чуждого осуществляется без участия третьей стороны и составляет измерение, отдельное от ''то же / другое'' (которое опирается на диалектически созданное целое). Выражаясь языком западной философии, другой (έτερον) и чуждый (ξένον) --- не одно и то же. Чуждость гостя --- в том числе Незнакомца из Элеи, появляющегося в \textit{Софисте} Платона,--- иного языка или культуры не сводится к тому что нечто кажется другим. Дерево и бетон отличаются друг от друга, но никто не станет говорить, что они чужды друг другу. Чуждость значит, что у \textit{я} (\textit{ipse}) есть сфера принадлежности, собственного бытия, и что это \textit{я} не то же что \textit{то же} (\textit{idem}), которое различимо лишь третьей стороной.

\section{Компромиссы модерна}

Ницше пишет: ''Важно избавиться от Вселенной, от единства ... разрушить Вселенную, разучиться уважать ее'' (KSA 12:317). И едва ли модерн серьезно относится к последствиям разграничения порядков. До сих пор мы стремимся к компромиссам через взаимное смешение \textit{я} и \textit{того же} и \textit{другого} и \textit{чуждого}. Двойственность оценки потенциала нового порядка видна и в конфликте модерна и постмодерна. Взять хотя бы \textit{я}, которое говорит ''я'' и так отделяет себя от целого. ''Я'' буквально значит исключение. Прямое соединение картезианского эго с эгоцентричностью было бы чрезмерным упрощением, ведь пренебрегло бы беспокойством, проистекающим из этого радикального самоанализа и побуждающим философа искать незыблемый фундамент. Поиск обречен на провал, но что им движет? Даже если ''\textit{Кто} я?'' слишком быстро соскальзывает к ''\textit{Что} я?'', этот вопрос, замыкая вопрошающего и вопрошаемое, уже содержит след ''Je est un autre'' Артюра Рембо.

Как бы то ни было, открытие \textit{я}, чья самость отделяет его от структуры отношений в природном и социальном мире, ослабевает, когда его превращают в лишь \textit{частность}. В конце концов мы бы столкнулись только с конфликтом \textit{индивидуализма} и \textit{холизма}, \textit{партикуляризма} и \textit{универсализма} --- конфликта все еще сохраняющегося но никогда не приводившего к значимым результатам. Если целое нельзя понять без частей, в которых оно выражено, а общие правила малопригодны без учета обстоятельств частностей, к которым применяются, то конфликт этот приводит к великому объединению: аристотеликов --- с кантианцами, герменевтов --- с универсальными прагматиками. Но гармония нарушается, если вспомнить, что элементы ''Я-Ты-Здесь-Сейчас'' не принадлежат к общим концептуальным категориям --- скорее, эти индексальные или окказиональные выражения являются указателями на место произнесения речи, место развертывания поля опыта, языка и действия, прежде чем оно само займет какое-либо место. Место, о котором говорится, не совпадает с местом произнесения речи, и то же касается момента времени. Даже как некто, говорящий ''я'', я не являюсь счетным элементом класса или частью целого. ''Я'' произношения, будучи явно отлично от ''я'' содержания произнесенного --- не исчислимая вещь, о которой можно было бы говорить во множественном числе, и то же касается ''ты''. В древнегреческой мысли ''я'' и ''мы'' не играют тематической роли за исключением понятия ''для нас'' (πρός ήμας). Взрывная сила я-речи также теряется, если сводить ''я'' к общей я-функции. Здесь нам нужно защитить Декарта от его трансцендентально-филосовских наследников.

То же затруднение касается и второго --- приписывающего всем порядкам неустранимую контингентность --- открытия. Если рассматривать \textit{способность быть другим} или \textit{способность действовать по-другому} как большую или меньшую степень \textit{произвольности}, то антиподом полагается \textit{необходимость}. Касательно претензий на легитимность, это ведет к еще одному конфликту --- релятивизма и универсализма, оружие которых сейчас уже столь же изношено. Здесь стоит еще раз напомнить: установления порядков с их легитимностью --- включая генеалогию истинного и ложного или добра и зла --- не верно ни относительно, ни абсолютно. Оно никак не верно, потому что существование бинарных стандартов не подчинено самим этим стандартам, если только их генезис опять не скрыт и соответствующая оппозиция не гипостазирована. У каждого порядка есть свое слепое пятно --- что-то неупорядоченное, что не просто нехватка. Это относится как к моральным, так и к когнитивным и эстетическим порядкам. Поэтому современные авторы так упорно борются с \textit{инципитами} своих романов: с самого начала и писатель, и читатель попадают в ловушку готового порядка.\footnote{См. \textit{Феноменология внимания} (\textit{Phänomenologie der Aufmerksamkeit}), глава 3 \textit{Неисчислимое} (\textit{Unerzählbares}).} Этот порядок постанавляет: ''есть порядки'', где ''есть'' не поддается никакому обоснованию, являясь базовым допущением всякой попытки обоснования --- нулевым уровнем. То, что сейчас называется постмодернизмом, может значить, среди прочего, что некоторые поля сражений утратили свое значение --- но не гарантирует, что проблемы, проистекающие из модерновой субъектности и рациональности, не возникнут вновь уже в другом месте.

\section{Парадоксы самоограничения}

Парадоксы нынче в моде. Парадокс --- παρά δόξαν --- буквально мнение, расходящиеся с преобладающим. Как и в случае других пара- (паразитизм, парафизика, паралогизм, парапсихология), понятие парадокса предполагает нормальное положение вещей. Покуда парадоксы отклоняются лишь от преобладающих мнений, к ним нет возражений. Совсем другое дело, когда предположение вступает в противоречие с собственными предпосылками, усиливая одно предположение за счет другого, и наоборот. Для сохранения предположений значение парадоксов обычно преуменьшают, как в случае с парадоксом лживого критянина. Конечно, различение объектного и метаязыка спасает высказывание --- его событие теперь не часть его содержания. Но такое методологическое решение основано на вышеупомянутом ограничении, приходящим извне и требующем различения между двумя уровнями речи. Анализ и решение проблемы находятся на уровне говоримого, что значит пренебрежение возможностью высказывания и само-высказывания --- сказанного, но не явленного в речи. Самореферентность есть, например, в ''Я обещаю вам, что ...'', из-за нее критянин вовлекает себя во все, что говорит. Кроме того, хорошо известное различие между плоскостью содержания и плоскостью отношений из теории коммуникации Пола Вацлавика не может быть прямо применено \textit{к} речи, потому что всегда присутствует \textit{в} ней.

Нечто подобное происходит и при самоограничении, когда ограничивающее себя является результатом ограничения, а не некой третьей стороной, отделяющей одно от другого, не будучи в них вовлеченной. Эксперт по грибам ни ядовит, ни съедобен, и судья не имеет отношения к деяниям обвиняемого и ущербу истца. Но что происходит, когда некто извиняется? Легко вспомнить случаи, когда мы сталкивались с чем-то странным, от чего \textit{сами} ограждаются --- без третьей стороны: разница между близостью и удаленностью как разница между ''здесь'' и ''тогда'', ''сейчас'' и ''давным-давно'', чередующиеся миры бодрствования и сна, сверхъестественные ''беседы с умершими'', переносящие из одного мира в другой через провал в памяти в романе Хуана Рульфо ''Педро Парамо''\footnote{Интересный взгляд на мотив чуждого в современной литературе и теории литературы см. в \textit{Мексика за гранью одиночества --- попытка межкультурного анализа} (\textit{Mexiko jenseits der Einsamkeit -- Versuch einer interkulturellen Analyse}, 1994).}, отношения между мужчиной и женщиной, ребенком и взрослым, социальная изоляция из-за классовой принадлежности, профессии или культуры, наконец, границы нормы, отделяющие здоровых от больных, невиновных от преступников, ортодоксов от еретиков, своих от чужих. Во всех этих случаях относящий себя к одной из сторон оказывается вовне. Такие перемены, как засыпание и пробуждение, болезнь и выздоровление, старость, обращение в другую веру значат не новый вид от третьего лица, а становление чем-то другим. Пруст показал, в какой степени пробуждение зависит от чисто телесной самоориентации, его \textit{Исследования} начинаются на грани сна и бодрствования. Такие переживания --- когда свое сталкивается с чуждым --- преимущественно переживания пересечения границ и, как станет видно, сопровождаются удвоением движения. Двойная природа ограничения, когда, с одной стороны, некто отделен от другого, а с другой --- определенные области разграничены, допускает дальнейшую дифференциацию.

Первый аспект касается самореферентности действия по ограничению. Действующая граница не просто находится или проводится. Иллюстративно проведенная на доске линия здесь вводит в заблуждения. Акт проведения границы, когда некто отделяет себя от другого, нельзя ни увидить, ни потрогать --- он схватывается только через свой след. Его можно сравнить с заключением контракта --- действием, не становящимся частью самого контракта, но являемым косвенно, в изменении обязанностей. Процесс проведения границы происходит в нулевой точке, не находящейся ни здесь, ни там. Действующая граница потому --- ни определенная вещь, ни ничто, ведь без нее не было бы ни того, ни другого --- ни меня, ни других. \textit{Самореферентность проведения границ заключается в их самоустранении.}

Второй аспект касается самореферентности, проявляющейся в \textit{само} самоограничения и приводящей к внутреннему и внешнему ограничению. Эта самость --- ни завуалированный субъект собственного акта ограничения, ни объективный результат чуждого акта ограничения: она как бы выскакивает в самом акте ограничения, проявляясь как полость, \textit{как внутреннее, отделяющее себя от внешнего} и так создающее \textit{предпочтение в этой разнице}. Формально это значит: то, \textit{что} отделяет себя, отмечается, тогда как то, \textit{от чего} оно отделяется, остается неотмеченным. В этом неустранимая асимметрия, без которой не было бы себя, способного принять точку зрения другого или третьей стороны. И вышеупомянутые примеры иллюстрируют эту --- присущую контрасту между своим и чуждым --- односторонность. ''Здесь'' и ''сейчас'', отделяя себя от расстояния, прошлого и будущего, относятся к разграничению, но они же представляют место, в котором это разграничение происходит. То же верно для всякого другого разграничения. Мы не имеем дело с нейтральной, бестелесной, бесполой и нестареющей третьей стороной, которая из своей нейтральности проводила бы различие между мужчиной и женщиной, взрослым и ребенком --- нет и транснациональной или транскультурной третьей стороны, могущей, например, настроить немцев против французов или европейцев против азиатов. Немцы, отграничивая себя от французов, могут быть теми, кто они есть, только благодаря этому разграничению. Выделяющий, обособляющий себя --- на одной стороне, чуждое, непонятное, от чего он обособляется --- на другой. \textit{Обращение к чуждому заключается в отстранении от него.} И удваиваясь и множась, эта асимметрия не приводит к симметрии. Симметрия требует подчиненности своего и чуждого общей точке зрения или общему правилу. ''Уравнивание неравных'', упомянутое Ницше в ''Об истине и лжи во вненравственном смысле'' предполагает равенство \textit{не}равного. Нормализация как избирательная форма упорядочения, начинаясь уже с ортоэстезии органов чувств, не может быть понята без влияния гетероэстезии, и то же относится к ортодоксии, ортологии и ортопраксии.\footnote{Эти идеи рассматриваются в моей работе \textit{Grenzen der Normalisierung} (\textit{Пределы нормализации}) (1998)}

Парадокс самореферентности усиливается еще и тем, что первый аспект самореферентного ограничения и второй --- заключающийся в ограничении себя внутрь и так ограничении чуждого вовне --- перетекают друг в друга в форме \textit{обращения к чуждому в обращении к себе}. От неизбежной самотематизации, ''я'' оказывается в системе упорядочивания, которую само же создает. ''Здесь'' тела можно определить, например, по красному маркеру на карте, указывающему местоположение пользователя. ''Сейчас'' имеет дату и отмечено в календаре. Тело (будет подробно рассмотренно в следующей главе) проявляет себя как живое тело, постоянно ускользающее от рефлексии сознания и сопротивляющееся нашим инициативам. То же происходит на социальном уровне и проявляется в межкорпоральности. ''Я'', без которого не было бы других, описывается Гуссерлем в \textit{Кризисе} как ''привелигированный член'' (\textit{Hua} VI:188; \textit{Crisis}, 185), но в то же время ''я'' обнаруживает себя обычным членом группы, которую помогает конституировать. ''Я'' --- это Другой, потому что чуждость рождается в своем. Обращение к чуждому в обращении к себе объясняет, почему никто не является просто тем, кто он есть, цепочка удвоения себя происходит при различных обстоятельствах в работах Гуссерля и Мерло-Понти, Фуко и Лумана. Эти удвоения не следует путать с отражениями рефлексивного самосознания, являющегося и субъектом, и объектом, не связаны они и с интерсубъективной диалектикой узнавания --- когда один субъект обнаруживает себя в другом субъекте. Телесное самоустранение, становясь очевидным в различии между функционированием и тематизацией, соответствует скорее различию между актом выражения и выраженным, между высказыванием и сказанным, ведь событие высказывания никогда не исчерпывается сказанным. Вообще, такого рода беспокойство, приводящее к постоянному удвоению ''я'' проистекает из того, что место, откуда проводятся границы, находится ни внутри, ни за пределами порядков, но сразу и внутри и снаружи них. Самоустранение значит опасность определенных моментов чужеродности в ''я'', в данном порядке. Двойная игра, как показывает Фуко, превращается в обманку, когда мы пытаемся уловить видение в виденном, высказывание в сказанном, мышление в мысли, вместо восприятия этих удвоений как части самой игры.

\section{Возможности и невозможности}

Тип мышления, подчиняющийся парадоксам самоограничивающегося --- и, значит, самопревосходящего --- порядка, порождает ряд моделей мышления, не вписывающихся в рамки традиционного мышления. \textit{Отклонения} от порядка играют важную роль не только в современной поэтике, но и в самоорганизации природы, они не указывают на существование нечта отклоняющегося, потому что отклонения проистекают скорее из самоизменения зрения или речи, которые ускользают от себя. Другой тип мышления --- мышление \textit{избытка}. Этот тип речи и действия, двигаясь вдоль границ их соответствующих порядков, отстает от самого себя, потому что привносит в игру больше возможностей, чем может использовать, и в то же время выходит за пределы себя, потому что касается невозможностей невидимого, неслышимого и немыслимого. Большее или меньшее количество смысла относится к формам неклассического искусства и указывает на избыток или нехватку смысла, чего не найти в классических паттернах исполнения желаний и приказов. Как ''не-фиксированное животное'', вынужденное изобретать и создавать порядки, человек раскрывает себя и как испытывающего нехватку, и как обладающего избытком: не может положиться даже на естественные потребности, не может сосредоточиться даже на требующих единого стандарта, регулятивных идеях. ''Есть порядки'' относится к тем порядкам, которые действительно подчиняются необходимым и ограничительным условиям, для которых, однако, не найти достаточных оснований. Порядки создают возможности и невозможности, но возможным не становится установление новых порядков. Здесь находится момент фактической безусловности среди опыта. Радикально чуждо то, что невозможно предвидеть ни субъективными ожиданиями, ни транссубъективными условиями возможности. Избыток опыта, выводя нас за рамки существующих порядков, в конечном счете приводит к пространственно-временному \textit{сдвигу}: самореференция никогда не возвращается к себе, место и время речи никогда не совпадает местом и временем \textit{в} речи. Самоустранение значит присутствие отсутствием, близость отдалением. Сдвиг начинается с рождения, рождение никогда не полностью собственное, ведь никогда не переживается активно, не выбирается. Причастность к собственному прошлому игриво выражает Лоренс Стерн, когда раз за разом откладывает рождение своего героя Тристрама Шенди, в то же время предвосхищая его во фразах вроде ''Я еще не родился'' (1950, стр.33). Карло Леви пишет об этом: ''Тристрам Шенди не хочет рождаться, потому что не хочет умирать'' (Italo Calvino, \textit{Lezioni Americane}, 1988, 46). Рождение касается меня, и все же не мой акт. Оно всегда преждевременно, новорожденный всегда запаздывает, задержка эта присуща всему новому, всему выходящему за рамки существующих мер. Каждое новое формирование проявляется как \textit{реформирование} существующих формаций, потому первая речь и первый поступок невозможны --- как и последняя речь и последний поступок. ''Абсолютное настоящее'', которое заключало бы в себе весь смысл --- один из фантазмов традиционных порядков, отрицающих свое происхождение.

Чуждое, выходящее за пределы игрового поля возможностей как нечто экстра-ординарное, можно потому рассматривать как нечто не-возможное --- невозможность здесь не онтологическая, эпистемологическая, практическая или логическая, а проживаемая.\footnote{Эта невозможность не ограничивается патологической стороной, а включает ее. Мой вклад см. в \textit{Gelebte Unmöglichkeit} под редакцией C.Rohde-Dachser и F.Wellendorf (2004) в \textit{Inszenierungen des Unmöglichen: Theorie und Therapie schwerer Persönlichkeitsstörungen}.} Дефис значит, что выходящее за рамки предписаний ведет не в потусторонний мир, а за пределы этого мира. Однако вопрос о возможностях и невозможностях опыта, речи и поступков демонстрирует двусмысленность, с которой наши рассуждения достигают своей крайней точки. В связи с этим возникает вопрос о связи между самоосознанием и обращением к чуждому. Это обращение можно истолковать двояко: как \textit{ограничение} своих возможностей или как \textit{сомнение в своем}.

Герменевты вроде Гадамера или прагматики вроде Рорти, полагаясь на принадлежность к традиции, занимают неопределенную промежуточную позицию: чуждое включается в обсуждение ровно настолько, насколько оно ему принадлежит, так что любое сомнение в своем обращается сомнением своего в своем, каждый остается сам по себе. И слабость рассудочной способности традиции, о которой говорит Джанни Ваттимо, мало что меняет. С другими теоретиками все однозначней. Первую интерпретацию обращению к чуждому мы находим в системной теории Никласа Лумана, он настойчив и в то же время гибок в размышлениях о том, как обращаться с границами, но всегда останавливается на определенной границе. Для Лумана обращение к чуждому значит претерпевшее сдвиг обращение к себе. Каждое наблюдение, что-то знача, и тем отделяя его от прочего, само по себе может стать предметом наблюдения для себя или другого. Единственное странное, что остается --- operatio pura наблюдателя, смутно напоминающее actus purus из аристотелевской теологии. Слепое пятно наблюдения не стирается, а перемещается с одной системной позиции на другую. В отличие от коммуникативного разума, всегда превращающего странное в нечто общее, системный оставляет место для чуждости, но чуждость эта --- \textit{исключенные, но при определенных условиях включаемые вновь, возможности}. Она лишь обозначает границы способностей --- личных, общих или оперативных. Таким образом, самореферентность первична, даже когда доступ к себе четко ограничен, и эта первичность превращается в свою полную противоположность, когда ''я'' подчиняется силе чуждого и так признает его своим безусловным хозяином. Похоже, обе крайности --- \textit{функционалистское непонимание} чуждого и \textit{фундаменталистская идеализация} чуждого --- стремятся усилить друг друга. Тогда мимолетной игре с возможностями была бы противопоставлена новая компактная реальность, как уже было показано в Какании Музиля. Моральные наблюдатели и моралисты обеспечивают друг друга уймой работы, стремясь нанести друг другу удар в спину: моральному наблюдателю напоминают о его наблюдательской морали, моралисту --- о его скрытой аморальности.

И это все? Допускают ли границы порядков только антитезу, только полемическую взаимодополняемость обращений к себе и к чуждому? Альтернативой, соответствующей второй форме обращения к чуждому, было бы возвращение к \textit{обращению к себе в обращении к чуждому}, с новой формой ответственности, позволяющей сочетать неизбежность требований с созданием наших собственных ответов. Как станет видно далее, эта не-возможность может также вернуть к требованиям чуждого, противоречащим нашим собственным идеям и желаниям и нарушающим общие правила. Сопротивление чуждого проистекает из отсутствия какого-либо эквивалента --- включая моральный,--- ускользающего из порядка. Для всякого порядка чуждое остается вторгающимся.

\chapter{Между претерпеванием и реакцией}

Если предположить, что чуждое нарушает границы любого порядка, то на что похож опыт такого нарушения? Не стоит ждать, что сочетание здравого смысла и правил, намеренных действий субъекта, руководящегося правилами, наряду с консенсуальным согллашением между субъектами, может противостоять вызовам чуждого. Рассматриваемая здесь альтернатива предстает в форме феноменологии, основанной на претерпевании и направленной на ответность, однако альтернатива эта уводит нас на задворки особого рода феноменологии и герменевтики, которые нацелены исключительно на интерпретацию смысла.

\section{В сфере смысла}

Подобно Фрейду, полагающему бессознательное основой психоанализа, можно положить интенциональность основой феноменологии. В точном определении интенциональность означает, что нечто проявляет себя как \textit{нечто},\footnote{Это ''как'' знакомо нам из других языков как \textit{qua, als, comme,} или \textit{come.}} что нечто подразумевается, дается, понимается или трактуется определенным образом: зеленая трава, пятно крови, вкус клубники, стол, Пегас, арабское число, любовное письмо, компьютерная программа, инфекция, нападение, террористический акт и т.д.. Формула \textit{нечто как нечто} означает, что нечто --- действительное, возможное или невозможное --- связано с чем-то другим (смысл, значение) и в то же время отделено от него. Реальность и чувство нельзя противопоставлять как характеристики или ценности. Подобно фуге, соединяющей несоединимое, миниатюра ''как'' обозначает разорванную связь. Мерло-Понти говорит о нулевой точке, стыке или складке, характеризующей возникновение смысла и Гештальта. ''Как'' не образует третьей сущности, перемещающейся между двумя исходными реальностями --- реальной и идеальной или физической и психической; скорее оно значит динамическую структуру, без которой буквально ничего проявляющего себя не было бы --- не было бы, значит, никого, кому оно является. Ничто не дается не будучи дано \textit{как нечто}, и никто не реагирует на него не реагируя \textit{как некто}. Опираясь на порождающий смысл аспект этого дифференцирующего процесса, я часто говорил об \textit{обозначающем различии}. По сравнению с основным характером этого дифференцирующего процесса, обращение к актам наделения смыслом предстает как специфическая интерпретация. Тот факт, что я замечаю оговорку, записку на двери или странный звук двигателя, еще не имеет характера приписываемого мной себе действия.

Если придерживаться подлинного дифференцирующего характера учения о значении, то обнаруживаются не только взаимосвязи между разными версиями феноменологии, но и связи с герменевтикой присутствия, интерпретацией традиций и текстов, семиотическими подходами, и, наконец --- что не менее важно --- с аналитической философией. Это делает идею смены парадигмы в феноменологии от философии сознания к семантике, предложенную Эрнстом Тугендхатом, гораздо менее правдоподобной. Несмотря на все разночтения между упомянутыми здесь альтернативами, на первый план выходит нечто общее, приводящее к философии смысла --- которая обычно заводит в тупик обычную эпистемологическую борьбу вроде таковой между реализмом и идеализмом, одновременно прекращая перетягивание каната между субъектом и объектом. Разделение на внешний и внутренний миры, введенное еще Джоном Локком и неоднократно после него повторявшееся, требует дополнительного третьего мира идей, поскольку оказывается конструкцией, покидающей почву опыта прежде ее обнаружения. Пронизанный интенциональностью, опыт не происходит ни внутри, ни снаружи. Так же и разделение эмпирических данных и общих идей оказывается подорванным. Как и в случае с формированием смысла и Гештальта, с самого начала опыт стремится к обобщению, не опираясь на данные в качестве основы. Присущие интенциональности ''как'' и ''каким образом'' подразумевают возможность повторения, и, следовательно, идеальность в генетическом и оперативном смысле, предшествующую всякой эйдетической или категориальной интуиции. Иначе говоря, интенциональный и дифференцированный опыт порождает множество горизонтальных и вертикальных опосредований, не полагаясь на заранее готовый разум или направляющего субъекта. Разум и субъект сами претерпевают генезис. Феноменология с концепцией горизонта смысла и опыта устанавливает совершенную пограничную концепцию и по-своему рассматривает вопрос об описании и преодолении границ без их устранения.

Обсуждаемое в первой главе мышление внутри и в пределах опыта также вытекает из базовой позиции, обозначенной как принцип феноменологии --- требуется лишь небольшой шаг, чтобы наделить \textit{нечто как нечто} коэффициентом контингентности. То что нечто предстает как нечто значит, что оно предстает \textit{только так и никак иначе}. С точки зрения теории Гештальта, каждый смысл, раскрывающийся как контекст ссылок --- привилегированный смысл: растение можно использовать как лекарственное или оставить как сорняк, нож можно использовать для еды или как оружие, денежный взнос может рассматриваться как пожертвование или как взятка, к иностранцу могут относиться как к просителю убежища или как к нелегалу. Как подчеркивал Ницше, привязанность к смыслу --- необратимая перспективность, но она же и потенциальная возможность конфликта. Есть смысл, но не какой-то отдельный смысл --- смысл развивается на фоне бессмыслицы, как показывает Мерло-Понти в своей \textit{Феноменологии восприятия} (Merleau-Ponty, 1945, 342; 2002, 345). Потому феноменология --- как и герменевтика --- повсюду приобретает индексальный характер. Происхождение системы Здесь-Сейчас, которую Карл Бюлер --- вдохновленный \textit{Логическими исследованиями} Гуссерля --- избрал в качестве центральной точки своей теории языка, не представляет собой комплекса принципов --- оно формирует поле, поле указаний, лежащее в основании символического поля с его расширениями из контекста и деконтекстуализациями. Помимо повседневных смысловых сочетаний, эта индексальность таит определенные исторические черты, свидетельствующие об эпохальных изменениях. Даже те события, которые мы обычно называем ''установлениями'' [\textit{Stiftungen}] --- рождение трагедии, демократии или геометрии --- имеют свои особенности, свои фазы инкубации и постановления. Их развитие не может быть подчинено телеологии, хотя близость понятия смысла к традиционному понятию цели или предназначения к этому манит.

Итак, я не вижу причин отказаться от этой плодотворной перспективы и занять другую. Поступать так было бы наивно: это подразумевает традиционное столкновение принципов или методов. Ориентация на смысл лежит в основе различных теорий значения, понимания и коммуникации и даже играет определенную роль в системно-теоретических подходах. Ее нельзя отрицать, не замалчивая. Но то, что нельзя отрицать, еще не окончательно оформлено и не гарантирует своей невосприимчивости к подрывным действиям, и что справедливо для смысловых формаций, то тем более справедливо для тех притязаний на достоверность, которые основаны на дискурсе и нормативно служат для фильтрации смысловых формаций. Речь никогда не о резких противопоставлениях или простых дополнениях: нас интересуют изменения, проявляющие свой устойчивый эффект подобно нюансам цветной живописи.


\section{Бездны претерпеваемого}

Сомнение в интенциональности должно вытекать из нее самой: как и в случае понимания или коммуникации, мы обречены на неудачу, хватки которой пытаемся избежать. Прямая конфронтация с ней означала бы возврат к лишенным смысла и норм данностям --- притворство, будто есть что-то незатронутое предписаниями смысла и норм. Предположение о таком ''чем-то'' рано или поздно попало бы в сети здравого смысла и было связано цепями аргументации. Следуя основной формуле, выбранной для этого анализа, мы должны понять --- как только внимательней рассмотрим возникновение смысла --- что есть ошибочные пути особого рода. То, что нечто проявляется как \textit{нечто}, не значит что это \textit{нечто}. Оно \textit{становится нечтом} через обретение смысла, и так становится высказываемым, доступным, повторяемым. Недостаточно предполагать происхождение смысла --- как если бы что-то просто выходило на свет из глубин опыта. Речь о достижении, порождении смысла, как говорят о генерации энергии. В этом гетерогенезе, проливающем также новый свет на гетерономию, вступает в игру нечто, что пока смысла не имеет, и вопрос в том, как понять этот \textit{гетеронт} как таковой.

Для прояснения позвольте привести отрывок из Роберта Музиля --- автора, не только близкого к Эрнсту Маху и вдохновлявшемуся Ницше, но и знакомого с ранними работами Гуссерля. В самом начале \textit{Человека без свойств} венская пара становится свидетелем события, которое мы обычно описываем \textit{как дорожно-транспортное происшествие} (так же как описываем некоторые метеорологические процессы и соответствующее им барометрическое давление, влажность и расположение солнца как \textit{великолепный августовский день}). Но вначале все не так. Все начинается с ''толпы'', с блокированного движения. Я слышу, как что-то тяжелое ''поворачивается в сторону, затем раздается звук резкого торможения, и то, что, по-видимому, является тяжелым грузовиком, застряло одним колесом на краю тротуара''. Рядом с ним стоит ''серый, как асфальт'' водитель и пытается объяснить, что произошло. Взгляды прохожих устремляются в ''глубину замкнутого пространства'', где жертва несчастного случая лежит как мертвая. Кто-то вызвал ''скорую'' --- ''экспертную, авторизованную помощь'', для которой экстренная помощь является нормальным явлением. Тем временем ''всезнайка'' объясняет какой-то женщине, что значит ''слишком большой тормозной путь''. Женщина благодарна за эти объяснения: ''для нее было достаточно того, что так это ужасное происшествие можно вписать в какую-то схему, превратить в техническую проблему, больше не касающуюся ее напрямую''. Это становится примером для ''социальных институтов'', которые так замечательно точны в своих услугах. ''У людей сложилось почти оправданное впечатление, что произошелшее было событием в рамках закона и правопорядка''. Комментарий разговорчивого мужчины и реакция его спутницы подтверждают это впечатление: '''Согласно американской статистике',-- отмечает мужчина,-- 'ежегодно на дорогах гибнет более ста девяноста тысяч человек и четыреста пятьдесят тысяч получают ранения'. 'Думаете, он мертв?' -- спрашивает его спутница, все еще питая неоправданное впечатление, что пережитое ею было чем-то исключительным'' (Musil 1978, 5-6).

Повседневное/неповседневное происшествие, которое здесь заканчивается, показывает, как трагическое событие превращается в статистическую норму, в ''упорядоченное'' событие: наделенное смыслом и соответствующее правилам, хотя все же оставляющее после себя субъективные ощущения в качестве прибавочного аффекта. Перед лицом растущей анонимизации кажущегося наиболее родным, это чувство напоминает личную реликвию: в конечном счете мы получаем ''переживания без переживающих'', предсказывает Музиль (1978, 150). Конечно, эту открывающую путешествие по Какании прелюдию можно было бы продолжить с помощью Фуко или Лумана, позволив, например, многофункциональной системе, называемой обществом, проявить себя со всеми ее органами чувств --- мы обнаружили бы различные коды и руководящие различия, будь то в медицинских тестах, юридических доказательствах или оценках действий, страховых агенствах, освещении инцидентов в средствах массовой информации --- когда уделяемое инциденту внимание зависит от общественной значимости жертвы и преступника,--- наконец, в литературном изложении, приобретающем эстетическую ценность в соответствии с разделением труда в этом производстве смысла, и все это не говоря уже о наших собственных феноменологических достижениях. Электрическая цепь смысла замыкается.

Безусловно, Гуссерль в \textit{Кризисе} высмеивает тех, кто приравнивает профессию феноменолога с таковой сапожника из пословицы, который всегда занят только своим делом\footnote{Let the cobbler stick to his last.} (Hua VI:140). Здесь можно было бы вновь процитировать Музиля, голосом главного героя смеющегося над ситуацией, когда ''птицеферма была создана ради возвышенных мыслей, называемых философией, теологией или литературой'' (Musil 1978, 358). Однако разделение профессий на специализации не должно немедленно приноситься в жертву какой-либо благородной цели: ни Гуссерль, ни Музиль этого не делают, еще меньше это делает Макс Вебер. Вполне возможно, что философ --- а следовательно, и феноменолог --- также сталкивается с определенными профессиональными требованиями, но он должен иметь способность и желание отразить профессионализацию образцовым образом. Музиль как литературный автор справляется с этой задачей: он не просто рассказывает о происшествии и вызывает тем сочувствие у читателя --- он ставит его перед фактом трансформации смысла значимого события. Есть нечто, что будучи хорошо подогнано и подготовлено по-другому, попадает в другие ''области смысла''. И все же, как возникает это ''незнамо-что'', что позволяет так легко представить себя в том или ином виде? Оно возникает как неожиданное событие, подобно corpus delicti ''ломает ряды'' и ''пересекает'' обычный ход событий, привлекая внимание, привлекая любопытных, разжигая желание увидеть, превращая наблюдателей в вуайерестов, возникает как то, что произошло с жертвой из-за ее ''беспечности''.

Несчастный случай, о котором говорится в романе --- прототип всего, что приходит нам на ум, что привлекает наше внимание, отталкивает, провоцирует, причиняет боль, заставляет думать и, в худшем случае, уничтожает нас как ''мыслящий тростник''. Здесь мы могли бы обратиться к феноменологии внимания, которая будет рассмотрена отдельно в одной из последующих глав. Что такое внимание? Обладает ли оно смыслом или это просто избирательная форма намерения? Так было бы только если бы внимание было сужено волюнтаристским образом или приравнивалось к лучу профектора, высвечивающего то, что всегда уже ожидает нашего взгляда в темноте. Однако привлечение внимания или осознания --- не те действия, которые можно совершать или от которых можно отказаться. Что-то может привлечь наше внимание как в виде легких, кажущихся безобидными, отвлекающих факторов, так и в форме сильных потрясений и опасностей --- как в описанном случае жертвы дорожного происшествия, которая следует своим намерениям, когда ее жестоко отрывают от них. Для обозначения этих --- не находящихся в нашем распоряжении, не ждущих разрешения или приказа, а просто происходящих с нами, одолевающих, волнующих, удивляющих нас, нападающих на нас --- событий, лучше подходит старый термин ''pathos'' (претерпевание), который, согласно высказыванию πάθει μάθος, говорит о познании через страдание, но не о познании страдания.\footnote{Старую доктрину катарсиса также необходимо бы снова подвергнуть сомнению --- вне зависимости от того, понимается ли катарсис как очищение аффектов или от аффектов. Процесс, понимаемый наполовину с медицинской, наполовину с моральной точки зрения, и ранее называвшийся ''очищением'', а сегодня --- ''проработкой'', слишком легко превращается в эмоционально управляемое упорядочивание претерпеваемого.} Такая перспектива освободит нас от нагромождения субъективных, личных эмоциональных состояний, проистекающих из ''дополнительной абстракции'' (\textit{Hua} VI:321; \textit{Crisis}, 226) и служащих для компенсации --- в форме психических данных --- утраченного при сведении мира к данным физическим. Но две вместе взятые абстракции не приводят к конкретизации: не хватает среды, в которой они могли бы ''срастись''.

Претерпевание не значит, что есть нечто воздействующее на нас, не значит оно и что нечто понимается и интерпретируется как нечто иное. И меньше, и больше: это понятие избегает альтернативы причинности и интенциональности во всех их традиционных формах.\footnote{Что интенциональность и причинность не разделяет пропасть, стало ясно мне довольно рано (см. \textit{Der Spielraum des Verhaltens}, 1980, глава 4). Но потребовалось некоторое время для достаточно точной дифференциации концепции эффекта.} Полагайся мы только на причинно-следственную связь, нам пришлось бы судить с точки зрения наблюдателя, подобно полицейскому, берущему образец крови, или врачу, ищущему микробы, борющемуся с риском заражения и определяющему причину смерти. Если же только на чистые процессы постижения, интерпретации и анализа --- это значило бы, что мы уже вышли на уровень осмысленности опыта, где его переломные моменты уже скрыты за интерпретацией смысла.

\end{document}
